%!TEX program = uplatex
\RequirePackage{plautopatch}

\documentclass[uplatex,dvipdfmx]{jlreq}

% 余白調整（1ページ収め）
\usepackage[top=18mm,bottom=20mm,left=20mm,right=20mm]{geometry}

\usepackage{amsmath,amssymb}
\usepackage{url}

\pagestyle{empty}

% 行間少し詰める
\renewcommand{\baselinestretch}{0.96}

%--------------------------------
% 講演マクロ
%--------------------------------

\newcommand{\talk}[3]{%
\noindent
#1\quad #2\par
\hspace*{2em}#3\par\smallskip
}

%--------------------------------
% タイトル情報
%--------------------------------

\newcommand{\EwMnumber}{第38回}
\newcommand{\EwMtitle}{幾何学と表現論\\[2mm]\large Kostant–関口対応をめぐって}
\newcommand{\EwMdate}{2006年12月8日（金）14:30～18:00 ～ 12月9日（土）10:30～17:00}
\newcommand{\EwMplace}{東京都文京区春日1-13-27 中央大学理工学部}

%--------------------------------
% タイトル表示
%--------------------------------

\newcommand{\EwMtitleblock}{

\vspace*{-5mm}

\begin{center}

{\Large ENCOUNTER with MATHEMATICS}

\vspace{2mm}

{\large 数学との遭遇}

\vspace{4mm}

{\Large \bfseries \EwMnumber\ \EwMtitle}

\vspace{4mm}

\EwMdate

\vspace{2mm}

於：\EwMplace

\end{center}

\vspace{2mm}

}

\begin{document}

\EwMtitleblock

%--------------------------------
% プログラム
%--------------------------------

\section*{プログラム}

\subsection*{12月8日（金）}

\talk
{14:30～15:20}
{ベキ零軌道の魅力}
{関口 次郎 氏（東京農工大・工）}

\talk
{15:40～16:40}
{超ケーラー多様体入門}
{中島 啓 氏（京大・理）}

\talk
{17:00～18:00}
{ベキ零軌道の関口対応}
{落合 啓之 氏（名大・多元数理）}

\subsection*{12月9日（土）}

\talk
{10:30～12:00}
{特性サイクルの理論とその実リー群の無限次元表現論への応用について I}
{竹内 潔 氏（筑波大・数学系）}

\talk
{14:00～15:30}
{インスタントンと巾零軌道}
{中島 啓 氏（京大・理）}

\talk
{16:00～17:00}
{特性サイクルの理論とその実リー群の無限次元表現論への応用について II}
{竹内 潔 氏（筑波大・数学系）}

\bigskip

17:10～　懇親会（ワイン・パーティー）

%--------------------------------
% 案内文
%--------------------------------

\bigskip

別紙の趣旨に沿った集会の第38回を以上のような予定で開催いたします。  
非専門家向けに入門的な講演をお願い致しました。  
多く方々のご参加をお待ちしております。  
講演者による講演内容へのご案内を別紙に添付いたしますのでご覧下さい。

%--------------------------------
% 連絡先
%--------------------------------

\section*{連絡先}

112-8551 東京都文京区春日1-13-27 中央大学理工学部数学教室 

tel : 03-3817-1745  

ENCOUNTER with MATHEMATICS : e-mail : encmath@math.chuo-u.ac.jp  
homepage : \url{http://www.math.chuo-u.ac.jp/ENCwMATH}  

三松 佳彦 : yoshi@math.chuo-u.ac.jp / 高倉 樹 : takakura@math.chuo-u.ac.jp

\end{document}
